\documentclass{article}
\usepackage{amsmath, amssymb, amsthm}

\setlength{\parindent}{0pt}

\newtheorem*{claim}{Claim}
\newtheorem*{lemma}{Lemma}

\begin{document}

\begin{lemma}
    In the reduced echelon form of a matrix, a column is a pivot column if and only if it is not a linear 
    combination of the columns to its left.
\end{lemma}

\begin{proof}
    Suppose the column is a pivot column. Every column to its left has a $0$ in the pivot's row, but the pivot 
    column has a $1$ there, so no combination of them can produce it.

    \medskip

    Now suppose the column is not a pivot column, and there are $m$ pivots to its left. By definition of 
    echelon form, the entries below row $m$ of the column are all $0$. The pivot columns to its left are 
    $\mathbf{e}_1, \dots, \mathbf{e}_m$, so the column is a linear combination of them.
\end{proof}

\begin{claim}
    The reduced echelon form of a matrix is unique, i.e. if a matrix is reduced to a reduced echelon form 
    by two different sequences of row operations, the two resulting matrices are equal.
\end{claim}

\begin{proof}
    Let $A$ be an arbitrary matrix, and suppose some sequence of row operations reduces $A$ to a reduced 
    echelon form $R$ and another reduces $A$ to a reduced echelon form $S$. Row operations preserve the 
    solution set of $A\mathbf{x} = \mathbf{0}$, so $R\mathbf{x} = \mathbf{0}$ and $S\mathbf{x} = \mathbf{0}$ 
    have the same solutions. We show that this forces $R = S$.

    \medskip

    Suppose column $j$ of $R$ is not a pivot column. By the lemma, it is a combination of the columns to its 
    left, which gives a solution $\mathbf{x}$ with a $1$ in position $j$, zeroes after it, and the negated 
    coefficients before it. The same $\mathbf{x}$ also solves $S\mathbf{x} = \mathbf{0}$, so column $j$ of $S$ 
    is a combination of the columns to its left, and hence not a pivot column. By symmetry, $R$ and $S$ have 
    the same non-pivot columns, and therefore the same pivot columns (a column is either non-pivot or pivot).

    \medskip

    It remains to show that corresponding columns of $R$ and $S$ are equal. The $k$th pivot column of each is 
    in the same position and equals $\mathbf{e}_k$. For a non-pivot column, we write it as a linear 
    combination of the pivot columns to its left only, and take the corresponding solution $\mathbf{x}$ as 
    before. Since the pivot columns of $R$ and $S$ agree and this $\mathbf{x}$ also solves 
    $S\mathbf{x} = \mathbf{0}$, the non-pivot columns of $R$ and $S$ are equal. Therefore corresponding columns 
    of $R$ and $S$ are equal, so $R = S$.

    \medskip

    Hence the reduced echelon form of a matrix is unique.
\end{proof}

\end{document}
